\section{第~\ref{sec:sensors}~章　机器的输入}

传感器的结构已被直接测量——通过单细胞电流记录、耳蜗振动测量和视网膜细胞计数；灵敏度极限和散粒噪声公式来自物理学，而传感器编码被调节到信息最大，则是有有力支撑的假说。

{\footnotesize
\setlength{\tabcolsep}{3pt}\renewcommand{\arraystretch}{1.15}
\begin{longtable}{>{\raggedright}p{0.5cm}>{\raggedright}p{4.0cm}>{\raggedright}p{5.6cm}>{\raggedright}p{2.3cm}>{\raggedright\arraybackslash}p{1.7cm}}
\toprule
序号 & 命题 & 实验及测量内容 & 文献 & 状态 \\
\midrule\endfirsthead
\toprule
序号 & 命题 & 实验及测量内容 & 文献 & 状态 \\
\midrule\endhead
1 & 传感器是换能器：受体产生电流，眼睛和耳朵感觉细胞的输出是送上\nt{GLUT}总线的谷氨酸；最前端的细胞不发放脉冲（图~\ref{fig:transducer}） & 龟的视杆和视锥细胞释放兴奋性氨基酸：它被一小片切下的膜上的谷氨酸通道捕获。大鼠耳蜗内毛细胞：神经纤维末梢的电流经谷氨酸AMPA受体流过，释放频率随细胞的平缓去极化上升到$150$\,Hz & \cite{CopenhagenJahr1989,GlowatzkiFuchs2002} & 已测量 \\
2 & 暗处的光传感器以约一皮安的电流响应单个光子 & 用吸引电极记录蟾蜍视杆外段：对暗弱闪光的响应是量子化的，单个事件$\approx1$\,pA（暗电流的$5\%$），离散度$0.2$\,pA，效率$0.5$。人报告单个光子到达的正确率高于随机水平 & \cite{Baylor1979,Tinsley2016} & 已测量 \\
3 & 声音传感器能察觉小于一纳米的位移 & 穆斯堡尔法，豚鼠耳蜗基底膜$16$--$19$\,kHz处：在听神经响应阈值处，膜的速度约为$0.04$\,mm/s，即位移$\approx0.35$\,nm（我们的换算）。测量的是基底膜，而不是传感器的纤毛束 & \cite{Sellick1982,Hudspeth1989} & 已测量 \\
4 & 暗处的精度受光子散粒噪声~\eqref{eq:photonnoise}限制；光弱时积累时间更长 & 公式来自泊松统计。在视杆中，暗光下的噪声与独立量子事件之和相符；在明亮背景上，对闪光的响应更小、更短。“零点几秒”这个数字在此没有单独检验 & \cite{Baylor1979} & 理论 \\
5 & 输入范围：光为十个数量级，声为十二个数量级，而脉冲频率不到三个数量级 & 声音：基底膜在特征频率上的响应在$40$--$80$\,dB区间内斜率低至$0.2$\,dB/dB，$100$\,dB以上仍可见压缩。数量级本身是标准估计，本章未给出处 & \cite{Ruggero1997} & 已测量 \\
6 & 传感器响应为式~\eqref{eq:nakarushton}，而$\sigma$跟随平均亮度，使曲线沿对数轴移动（图~\ref{fig:adaptation}） & 公式最早拟合于鱼视网膜的S电位。龟的视锥细胞：响应服从$V=I V_m/(I+\sigma)$，覆盖$\sim3.5$个数量级的亮度；背景持续$2$--$3$分钟后，曲线水平移动而宽度不变。与除以总活动的联系见综述 & \cite{NakaRushton1966,NormannPerlman1979,Carandini2012} & 已测量 \\
7 & 传感器传送变化，其编码被调节到每个脉冲携带最多信息（命题~\ref{prop:sensors}） & 这一原则由理论提出。最有力的证据：苍蝇第一层神经元的“对比度$\to$响应”特性与自然场景对比度的累积分布一致——这样可获得最大信息 & \cite{Barlow1961,Laughlin1981} & 假说 \\
8 & 接通传感器和断开传感器构成双线编码；事件相机在硅片上复现了它 & 实验综述：视网膜的接通通道和断开通道在第一个突触处就已分开，并可分别关闭。相机：$128\times128$像素，无帧，动态范围$120$\,dB，延迟$15$\,\textmu s & \cite{Schiller1992,Lichtsteiner2008,Gallego2022} & 已测量 \\
9 & 眼睛有$\sim10^8$个光传感器和$\sim10^6$个输出神经元：压缩$\sim100$倍 & 在完整的人视网膜标本上计数：平均$92\cdot10^6$个视杆和$4.6\cdot10^6$个视锥；不同眼睛的输出（神经节）细胞为$0.7\cdot10^6$到$1.5\cdot10^6$个 & \cite{Curcio1990,CurcioAllen1990} & 已测量 \\
10 & 小鼠眼睛的输出通道超过三十种 & 小鼠：对$11\,000$多个输出细胞做双光子钙成像并对响应聚类——明显多于$30$种功能类型。人的数目尚未测量 & \cite{Baden2016} & 已测量 \\
11 & 豚鼠眼睛传送$\sim875\,000$\,bit/s；换算到人眼约为$\sim10^7$\,bit/s & 豚鼠，多电极阵列，自然场景中的运动：每个细胞$6$--$13$\,bit/s，$\sim10^5$个输出细胞共$\sim875\,000$\,bit/s。人（$\sim10^6$个细胞）的$\sim10^7$\,bit/s是作者的换算 & \cite{Koch2006} & 已测量 \\
12 & 耳朵是带通滤波器组，频率图近似对数（图~\ref{fig:filterbank}） & 基底膜最大振动的位置取决于频率；“频率—位置”函数近似指数型，用同一形式描述人和活体动物的耳蜗 & \cite{Greenwood1990,Robles2001} & 已测量 \\
13 & 谐振器是主动的：一部分传感器把微弱振动的振幅放大数千倍（$66$--$76$\,dB），响度增大时增益下降 & 激光测振，南美栗鼠耳蜗底部：弱纯音在特征频率上的增益相对镫骨为$66$--$76$\,dB；死亡使灵敏度降低$60$--$81$\,dB（振幅$10^3$--$10^4$倍）并消除压缩。动力来源是外毛细胞 & \cite{Ruggero1997,Robles2001} & 已测量 \\
14 & 耳朵的放大器工作在自激边缘 & 计算：在霍普夫分岔点附近，范围压缩、尖锐调谐和组合音不可避免——这正是听觉所有已知的非线性。耳蜗确实如此调节，尚无直接证明 & \cite{Eguiluz2000} & 假说 \\
15 & 低频时脉冲与声音同相（豚鼠的相位锁定从$\sim600$\,Hz开始减弱，直到$\sim3.5$\,kHz仍可察觉），并据此确定方向；高频时只剩位置编码 & 豚鼠听神经：相位锁定在约$600$\,Hz开始下降，$3.5$\,kHz以上消失；猫和猴的界限约高一个倍频程。双耳到达时间差见综述 & \cite{PalmerRussell1986,Grothe2010} & 已测量 \\
16 & 其他传感器（表~\ref{tab:sensors}）：嗅觉有$\sim400$种受体，每个神经元一种，组合编码；味觉部分经由ATP和血清素；触觉有快适应和慢适应传感器；冷和热各自独立 & 小鼠：一种受体识别多种气味，一种气味激活多种受体，不同气味对应不同集合。没有ATP受体，味觉神经不响应；味蕾释放血清素。人手皮肤单根纤维记录：按适应速度分为四类。冷通道不同于热通道 & \cite{Mombaerts2004,Malnic1999,Finger2005,Huang2005,JohanssonVallbo1979,McKemy2002} & 已测量 \\
\bottomrule
\end{longtable}}
